\documentclass{article} % 必须退回基础 article

% ==================== 导言区 (Preamble) ====================
\usepackage{booktabs} 
\usepackage{amssymb}  
\usepackage{amsmath}  
\usepackage{microtype} 
\usepackage{CJKutf8}    % 强制引入 pdfLaTeX 专用的 UTF-8 中文宏包

% ==================== 正文区 (Document Body) ====================
\begin{document}
\begin{CJK*}{UTF8}{gbsn} % 开启中文环境，使用简体宋体 (gbsn)

\subsection*{主实验}

本实验旨在观察当真实交通波动增强时，预测驱动的重新设计（forecast-driven redesign）是否仍然有效。
其中，DCRNN 和 MTGNN 模型使用原始经验误差（empirical error），即 $\alpha = 1.0$。
实验参数配置与对照如表~\ref{tab:main_experiment} 所示。

\begin{table}[htbp]
\centering
\small 
\caption{主实验设置与对照}
\label{tab:main_experiment}
\begin{tabular}{lcccc}
\toprule
$CV_T$ & No forecast & DCRNN noisy & MTGNN noisy & Perfect forecast \\
\midrule
0.05 & \checkmark & \checkmark & \checkmark & \checkmark \\
0.25 & \checkmark & \checkmark & \checkmark & \checkmark \\
0.50 & \checkmark & \checkmark & \checkmark & \checkmark \\
1.00 & \checkmark & \checkmark & \checkmark & \checkmark \\
\bottomrule
\end{tabular}
\end{table}

\subsection*{敏感性实验}

本实验旨在分析优化结果对预测误差的敏感性，重点回答：当预测误差逐渐增大时，优化结果从什么时候开始明显变差。
实验通过固定 $CV_T$（例如选择 $CV_T = 0.50$ 或 $1.00$），并改变预测误差缩放因子 $\alpha$ 的值来实现，$\alpha$ 的具体物理含义见表~\ref{tab:sensitivity_alpha}。

\begin{table}[htbp]
\centering
\caption{预测误差缩放因子 $\alpha$ 的含义}
\label{tab:sensitivity_alpha}
\begin{tabular}{cl}
\toprule
$\alpha$ & 含义 \\
\midrule
0.0 & Perfect forecast \\
0.5 & Low error \\
1.0 & Empirical model error \\
1.5 & High error \\
2.0 & Severe error \\
\bottomrule
\end{tabular}
\end{table}

\end{CJK*} % 结束中文环境
\end{document}